% !TEX program = xelatex
% 공통수학2 · Ⅰ. 도형의 방정식 · 3. 도형의 이동 · 02 대칭이동 · 2/2차시
% 학생용 활동지 (답안 없음)
\documentclass[11pt,a4paper]{article}
\usepackage{kotex}
\usepackage{iftex}
\ifPDFTeX\else
  \setmainhangulfont{Noto Serif CJK KR}
  \setsanshangulfont{Noto Sans CJK KR}
\fi
\usepackage[margin=2.1cm]{geometry}
\usepackage{parskip}
\usepackage{enumitem}
\usepackage{array}
\usepackage{tabularx}
\usepackage{xcolor}
\usepackage{amssymb}
\usepackage{tikz}
\usepackage[most]{tcolorbox}
\usepackage{fancyhdr}
\usepackage{titlesec}

\definecolor{goldc}{HTML}{B07C22}
\definecolor{goldstrong}{HTML}{8A5F16}
\definecolor{indigoc}{HTML}{2B3B57}
\definecolor{indigosoft}{HTML}{6E82A6}
\definecolor{linec}{HTML}{D6DACB}
\definecolor{surface2}{HTML}{E4E8DC}
\definecolor{writeline}{HTML}{C7CCBA}
\definecolor{inksoft}{HTML}{52604F}

\titleformat{\section}{\normalfont\sffamily\bfseries\large\color{indigoc}}{\thesection}{0.6em}{}
\titlespacing{\section}{0pt}{14pt}{6pt}
\newtcolorbox{goalbox}{enhanced, breakable, colback=white, colframe=goldc, boxrule=0.5pt, leftrule=3pt, arc=2pt, left=10pt, right=10pt, top=8pt, bottom=8pt}
\newtcolorbox{sheetbox}{enhanced, breakable, colback=white, colframe=linec, boxrule=0.6pt, arc=3pt, left=11pt, right=11pt, top=9pt, bottom=9pt}
\newcommand{\writespace}[1]{%
  \par\nobreak\vspace{3pt}
  \foreach \n in {1,...,#1} {%
    \noindent\textcolor{writeline}{\rule{\linewidth}{0.4pt}}\par\vspace{0.62cm}%
  }
}
\newcommand{\fillblank}[1]{\makebox[#1]{\hrulefill}}
\pagestyle{fancy}\fancyhf{}\renewcommand{\headrulewidth}{0pt}
\fancyfoot[L]{\footnotesize\color{inksoft} 공통수학2 · 2026학년도 2학기 · 지평선고등학교}
\fancyfoot[R]{\footnotesize\color{inksoft} 다음 시간 --- 중단원·대단원 마무리}

\begin{document}

\noindent
\begin{minipage}[b]{0.62\linewidth}
  {\sffamily\footnotesize\color{indigosoft} 공통수학2 · Ⅰ. 도형의 방정식 · 3. 도형의 이동}\\[2pt]
  {\sffamily\bfseries\Large 02 대칭이동 \textperiodcentered\ 16차시}
\end{minipage}\hfill
\begin{minipage}[b]{0.36\linewidth}
  \raggedleft\small\color{inksoft}
  반\ \fillblank{1.6cm} \quad 번호\ \fillblank{1.6cm}\\[4pt]
  이름\ \fillblank{3.6cm}
\end{minipage}

\vspace{4pt}
\noindent{\color{goldc}\rule{\linewidth}{2.2pt}}
\vspace{10pt}

\begin{goalbox}
\textbf{오늘의 목표}\\
도형의 방정식을 $x$축, $y$축, 원점, 직선 $y=x$에 대하여 대칭이동하고, 실생활 문제(반사)에 적용할 수 있다.
\vspace{4pt}
\small\color{inksoft}
$\square$ 자신 있음 \quad $\square$ 조금 헷갈림 \quad $\square$ 처음 봄 \hfill (수업 전 나의 상태에 표시)
\end{goalbox}

\section{개념 정리 --- 도형의 대칭이동}

\begin{sheetbox}
방정식 $f(x,y)=0$이 나타내는 도형을
\begin{enumerate}[leftmargin=1.6em, itemsep=3pt]
  \item $x$축에 대하여 대칭이동한 도형의 방정식은 $f(\ \fillblank{1.4cm}\ ,\ \ \fillblank{1.4cm}\ )=0$
  \item $y$축에 대하여 대칭이동한 도형의 방정식은 $f(\ \fillblank{1.4cm}\ ,\ \ \fillblank{1.4cm}\ )=0$
  \item 원점에 대하여 대칭이동한 도형의 방정식은 $f(\ \fillblank{1.4cm}\ ,\ \ \fillblank{1.4cm}\ )=0$
  \item 직선 $y=x$에 대하여 대칭이동한 도형의 방정식은 $f(\ \fillblank{1.4cm}\ ,\ \ \fillblank{1.4cm}\ )=0$
\end{enumerate}

\vspace{8pt}
\textbf{예제} --- 원 $(x-4)^2+(y+3)^2=4$를 $x$축, $y$축, 원점에 대하여 각각 대칭이동한 원의 방정식을 구해 보자.
\writespace{4}
\end{sheetbox}

\section{연습 문제}

\begin{sheetbox}
\textbf{1.} 다음 방정식이 나타내는 도형을 $x$축, $y$축, 원점에 대하여 대칭이동한 도형의 방정식을 구하시오.\\
\hspace*{1.6em}⑴ $2x-3y-1=0$ \qquad ⑵ $x^2+y^2-4x+10y-7=0$
\writespace{5}

\vspace{6pt}
\textbf{2.} 원 $(x-3)^2+(y+2)^2=5$를 $x$축의 방향으로 $2$만큼, $y$축의 방향으로 $1$만큼 평행이동한 다음 $x$축에 대하여 대칭이동한 원의 방정식을 구하시오.
\writespace{3}

\vspace{6pt}
\textbf{3.} 원 $x^2+(y-3)^2=9$를 직선 $y=x$에 대하여 대칭이동한 원의 방정식을 구하시오.
\writespace{2}
\end{sheetbox}

\section{실생활 탐구 --- 당구대 위의 반사}

\begin{sheetbox}
가로 $80$, 세로 $40$인 직사각형 당구대(두 변이 좌표축 위에 있음)에서, 점 P$(10,30)$의 노란 공을 쳐서 $y$축, $x$축에 차례로 부딪히게 하여 점 Q$(70,10)$의 빨간 공을 맞히려고 한다. 입사각과 반사각이 같다고 할 때, $y$축 위의 반사점 A와 $x$축 위의 반사점 B의 좌표를 구해 보자.

\vspace{4pt}
{\footnotesize\color{inksoft} 힌트: Q를 $x$축에 대하여, P를 $y$축에 대하여 각각 대칭이동한 뒤, 두 대칭점을 잇는 직선을 그리면 이 직선이 A, B를 지난다.}
\writespace{6}
\end{sheetbox}

\section{사고의 가시화 --- Connect · Extend · Challenge}

\noindent{\footnotesize\color{inksoft} `대칭이동' 전체(15$\sim$16차시)를 떠올리며 아래 세 칸을 채워 보자.}\\[6pt]
\begin{tabularx}{\linewidth}{@{}XXX@{}}
\textbf{\footnotesize\color{indigoc} Connect --- 무엇과 연결되나?} &
\textbf{\footnotesize\color{indigoc} Extend --- 어디까지 확장됐나?} &
\textbf{\footnotesize\color{indigoc} Challenge --- 아직 궁금한 것은?} \\[4pt]
\end{tabularx}
\noindent
\begin{minipage}[t]{0.32\linewidth}\writespace{3}\end{minipage}\hfill
\begin{minipage}[t]{0.32\linewidth}\writespace{3}\end{minipage}\hfill
\begin{minipage}[t]{0.32\linewidth}\writespace{3}\end{minipage}

\section{마무리 성찰}

\begin{sheetbox}
\begin{minipage}[t]{0.48\linewidth}
{\footnotesize\color{inksoft} `도형의 이동' 소단원을 마치며 한 문장으로}
\writespace{2}
\end{minipage}\hfill
\begin{minipage}[t]{0.48\linewidth}
{\footnotesize\color{inksoft} 감사 일기 / 유념 한 줄}
\writespace{2}
\end{minipage}
\end{sheetbox}

\end{document}
